% 自绘：g(a)=-tanh(a) 的特征投影；虚线仅表示光滑失效后的轨道投影。
\begin{figure}[htbp]
\centering
\begin{tikzpicture}[x=2.3cm,y=2.2cm,>=stealth]
 \draw[->] (-1.35,0)--(1.35,0) node[right]{\(x\)};
 \draw[->] (-1.3,-0.05)--(-1.3,1.6) node[above]{\(t\)};
 \draw[gray,densely dotted] (-1.25,1)--(1.25,1);
 \node[left] at (-1.3,1){\(1\)};
 \foreach \a/\v in {-1/-0.761594,-0.75/-0.635149,-0.5/-0.462117,-0.25/-0.244919,0/0,0.25/0.244919,0.5/0.462117,0.75/0.635149,1/0.761594}{
   \draw[AnalysisBlue,thick] (\a,0)--({\a-\v},1);
   \draw[AnalysisBlue,dashed] ({\a-\v},1)--({\a-1.45*\v},1.45);
 }
 \foreach \x in {-1,0,1}{\draw (\x,0.025)--(\x,-0.025) node[below]{\(\x\)};}
 \node[anchor=west,align=left] at (0.68,1.28){\small 轨道投影};
 \node[anchor=west] at (0.68,0.72){\small 光滑阶段};
\end{tikzpicture}
\caption[特征线的压缩与投影交叉]{特征线的压缩与投影交叉。初值为 \(g(a)=-\tanh a\)，轨道投影为 \(x=a-t\tanh a\)。在 \(t=1\) 首次发生梯度爆破；虚线延长只展示投影的交叉，不表示单值经典解的延续。}
\label{b2:fig:07-characteristic-compression}
\end{figure}
